% from ega4-18-fr.tex
\oldpage[IV-4]{217}
récurrence sur $b$ pour $b\geq1$ (le cas $b=1$ étant trivial) pour achever de prouver (19.7.1.16) et par suite aussi $(**_{d+1,h+1})$ pour $h\leq m\setmin d$.

\begin{env}[19.7.1.18]
\label{IV.19.7.1.18}
\emph{Démonstration de} $(**_{d+1,h+1})$; \emph{Deuxième cas}: $h>m\setmin d$.
\end{env}

Nous allons tout d'abord démontrer que, pour \emph{tout} $h\geq0$, on a
\begin{env}[19.7.1.19]
\label{IV.19.7.1.19}
\[
  R\cap(\mathfrak K^{d+1}N+\mathfrak K^d\mathfrak L^hN)\subset\mathfrak L^h(R\cap\mathfrak K^dN)+\mathfrak L^{h+1}\mathfrak K^dN+\mathfrak K^{d+1}N.
\]
\end{env}

Considérons pour cela un élément $z$ du premier membre de (19.7.1.19). Notons que l'on a $\mathfrak K^{d+1}N+\mathfrak K^d\mathfrak L^hN=\mathfrak K^{d+1}N+\mathfrak K^d\mathfrak J^hN$ puisque $\mathfrak L=\mathfrak K+\mathfrak J$; nous allons considérer la classe $\bar z$ de $z$ dans $\operatorname{gr}_{\mathfrak J}^h(\operatorname{gr}_{\mathfrak K}^d(N))$, qui, en vertu de (19.5.5) et de l'hypothèse sur $\mathbf f$, s'identifie à un sous-$(A/\mathfrak L)$-module de $\operatorname{gr}_{\mathfrak L}^{h+d}(N)$. Du fait que $z\in R$, nous allons montrer que l'on a même (avec l'identification précédente)
\begin{env}[19.7.1.20]
\label{IV.19.7.1.20}
\[
  \bar z\in\operatorname{gr}_{\mathfrak L}^{h+d}(R,N).
\]
\end{env}

En effet, on a noté, dans la démonstration de (19.7.1.3) que, pour tout idéal premier $\mathfrak r$ de $A$ tel que $\mathfrak p=\mathfrak r/\mathfrak K\in\operatorname{Ass}_{A/\mathfrak K}(A/\mathfrak L)$, l'application $\psi_{R_{\mathfrak r}}'$ est bijective, et par suite la relation (19.7.1.6), où l'on remplace tous les modules par leurs localisés en $\mathfrak r$, est vraie. Comme $z$ appartient à $R\cap(\mathfrak K^dN+\mathfrak L^{d+1}N)$, son image $z/1$ dans $R_{\mathfrak r}$ appartient à $(R_{\mathfrak r}\cap\mathfrak K_{\mathfrak r}^dN_{\mathfrak r})+\mathfrak L_{\mathfrak r}^{d+1}N_{\mathfrak r}$, donc l'image $\bar z/1$ de $\bar z$ dans $\operatorname{gr}_{\mathfrak L_{\mathfrak r}}^\bullet(N_{\mathfrak r})$ appartient à $\operatorname{gr}_{\mathfrak L_{\mathfrak r}}^\bullet(R_{\mathfrak r},N_{\mathfrak r})=(\operatorname{gr}_{\mathfrak L}^\bullet(R,N))_{\mathfrak r}=(\operatorname{gr}_{\mathfrak L}^\bullet(R,N))_{\mathfrak q}$, où $\mathfrak q=\mathfrak r/\mathfrak L$. En d'autres termes l'image de $\bar z$ par l'application canonique $(\operatorname{gr}_{\mathfrak L}^\bullet(N))_{\mathfrak q}\longrightarrow(\operatorname{gr}_{\mathfrak L}^\bullet(M))_{\mathfrak q}$ est nulle pour tout $\mathfrak q\in\operatorname{Ass}_{A/\mathfrak L}(A/\mathfrak L)$. Comme on a vu dans la démonstration de (19.7.1.3) que $\operatorname{Ass}_{A/\mathfrak L}(\operatorname{gr}_{\mathfrak L}^\bullet(M))\subset\operatorname{Ass}_{A/\mathfrak L}(A/\mathfrak L)$, on en déduit bien (Bourbaki, \emph{Alg. comm.}, chap. IV, § 1, n\textsuperscript{o} 1, cor. 1 de la prop. 2 et n\textsuperscript{o} 3, prop. 7) que l'image de $\bar z$ par l'application canonique $\operatorname{gr}_{\mathfrak L}^\bullet(N)\to\operatorname{gr}_{\mathfrak L}^\bullet(M)$ est nulle, c'est-à-dire que l'on a la relation (19.7.1.20).

Cela étant, on peut écrire par définition
\begin{env}[19.7.1.21]
\label{IV.19.7.1.21}
\[
  z\equiv\sum_{|\mathbf p|=h}c_{\mathbf p}\mathbf f^{\mathbf p}\ \mathrm{mod.}\,\mathfrak K^{d+1}N
\]
\end{env}
où on pose comme d'ordinaire $\mathbf p=(p_1,\ldots,p_n)$, $\mathbf f^{\mathbf p}=f_1^{p_1}f_2^{p_2}\ldots f_n^{p_n}$, $|\mathbf p|=p_1+\ldots+p_n$ et où $c_{\mathbf p}\in\mathfrak K^dN$. En outre, comme $\psi_N$ identifie $\operatorname{gr}_{\mathfrak L}^\bullet(N)$ à $(\operatorname{gr}_{\mathfrak K}(N)\otimes_{A/\mathfrak K}(A/\mathfrak L))[T_1,\ldots,T_n]$, les $c_{\mathbf p}$ sont \emph{déterminés} mod. $\mathfrak K^d\mathfrak LN$, et si $\bar c_{\mathbf p}$ est la classe de $c_{\mathbf p}$ mod. $\mathfrak K^d\mathfrak LN$, on a, par l'identification précédente,
\begin{env}[19.7.1.22]
\label{IV.19.7.1.22}
\[
  \bar z=\sum_{|\mathbf p|=h}\bar c_{\mathbf p}\mathbf T^{\mathbf p}.
\]
\end{env}
Utilisant (19.7.1.3), on en déduit que l'on a nécessairement, pour tout $\mathbf p$ tel que $|\mathbf p|=h$ ($\bar c_{\mathbf p}$ étant cette fois identifié par $\psi_N$ à un élément de $\operatorname{gr}_{\mathfrak L}^d(N)$)
\begin{env}[19.7.1.23]
\label{IV.19.7.1.23}
\[
  \bar c_{\mathbf p}\in\operatorname{gr}_{\mathfrak L}^d(R,N).
\]
\end{env}
Mais puisque $d<m$, l'hypothèse que $(*_d)$ est vraie implique que $\bar c_{\mathbf p}$ appartient à l'image de $\psi_R'$, donc qu'on peut supposer que $c_{\mathbf p}\in(R\cap\mathfrak K^dN)+\mathfrak K^d\mathfrak LN$. La relation (19.7.1.21) donne alors
\[
  z\in\mathfrak J^h(R\cap\mathfrak K^dN)+\mathfrak J^h\mathfrak K^d\mathfrak LN+\mathfrak K^{d+1}N
\]
et comme $\mathfrak J\subset\mathfrak L$, cela prouve (19.7.1.19).

